%% Capítulo 3 — Trabalho Relacionado
\chapter{Trabalho Relacionado}
\label{chap:trabalho_relacionado}

A Tabela~\ref{tab:comparacao_iam} compara as principais soluções \gls{iam} consideradas.

\begin{table}[h!]
\centering
\caption{Comparação de soluções IAM}
\label{tab:comparacao_iam}
\begin{tabularx}{\textwidth}{|l|l|X|X|}
\hline
\textbf{Solução} & \textbf{Modelo} & \textbf{Pontos fortes} & \textbf{Limitações} \\ \hline
Auth0 & SaaS & Fácil integração, documentação extensa & Custo elevado, dependência de terceiro \\ \hline
Azure AD B2C & Cloud & Integração Microsoft, escala empresarial & Configuração complexa, custo por login \\ \hline
Okta & SaaS & UX excelente, \gls{mfa} avançada & Licenciamento por utilizador \\ \hline
\textbf{Keycloak} & Open Source & \textbf{Gratuito, self-hosted, Admin API completa} & Curva de aprendizagem \\ \hline
\end{tabularx}
\end{table}

O \textbf{Keycloak~26} \cite{KeycloakDocs} foi escolhido por: custo zero; execução local via Docker; \texttt{realm-export.json} versionável; suporte nativo a \gls{mfa} \gls{totp}; Admin \gls{rest} \gls{api} para automação \gls{jml}; e revogação de sessões com efeito imediato via \texttt{DELETE /users/\{id\}/sessions}.

Relativamente à validação de tokens, optou-se por \textbf{validação local via \gls{jwks}} (em vez de \textit{token introspection}) para eliminar latência e dependência de disponibilidade do \gls{idp} em cada pedido, compensando a limitação de revogação imediata através da revogação explícita de sessões nos fluxos \gls{jml}.
